% ch-a4-05.tex —— 第 4 篇（Emerton–Gee–Hellmann）续片，底本 PDF 228–239 = 印刷页 215–226
\ysrc{215}
其中右端的张量积由某个（具体的）映射 $\mathcal{R}_F\to\widehat{F}_\infty$ 诱导。该映射 $\mathcal{R}_F\to\widehat{F}_\infty$ 经过 $F_\infty\subset\widehat{F}_\infty$ 分解，而模 $\widehat{D}_{\mathrm{Sen}}(D_A)$ 有一个去完备化（decompleted）版本
\[
D_{\mathrm{Sen}}(D_A)=D_A\otimes_{\mathcal{R}_{F,A}}(A\otimes_{\Q_p}F_\infty),
\]
它是有限投射 $A\otimes_{\Q_p}F_\infty$-模，带有 $1$ 维 $p$-进 Lie 群 $\Gamma$ 的连续作用。事实上，只要 $n$ 相对于 $D_A$ 足够大，$D_{\mathrm{Sen}}(D_A)$ 甚至可以在 $A\otimes_{\Q_p}F_n$ 上定义：即存在某个 $n\gg0$ 与（唯一的）$A\otimes_{\Q_p}F_n$-模 $D_{\mathrm{Sen},n}(D_A)$，带有半线性连续 $\Gamma$-作用，使得
\[
D_{\mathrm{Sen}}(D_A)=D_{\mathrm{Sen},n}(D_A)\otimes_{F_n}F_\infty.
\]
在去完备化模 $D_{\mathrm{Sen}}(D_A)$ 与 $D_{\mathrm{Sen},n}(D_A)$ 上，$\Gamma$-作用是局部解析的；但提醒读者，它并不光滑，也就是说，它不是下降到 $A\otimes_{\Q_p}F$-模的下降数据（descent datum）。

由于 $\Gamma_n\subset\Gamma$ 在 $D_{\mathrm{Sen},n}(D_A)$ 上的作用是 $A\otimes_{\Q_p}F_n$-线性的且局部解析，可对该作用求导，从而定义态射
\[
\operatorname{Lie}\Gamma=\operatorname{Lie}\Gamma_n\longrightarrow\operatorname{Lie}\bigl(\mathrm{GL}(D_{\mathrm{Sen},n}(D_A))\bigr)=\End_{A\otimes_{\Q_p}F_n}(D_{\mathrm{Sen},n}(D_A)).
\]
分圆特征标给出平凡化 $\Z_p\cong\operatorname{Lie}\Gamma$，因此 $1\in\Z_p$ 的像（在把系数由 $F_n$ 扩回 $F_\infty$ 之后）定义了 $A\otimes_{\Q_p}F_\infty$-线性自同态
\[
\Theta: D_{\mathrm{Sen}}(D_A)\longrightarrow D_{\mathrm{Sen}}(D_A),
\]
它与上述讨论中 $n$ 的选取无关。该自同态与 $\Gamma$-作用交换，因此其特征多项式的系数取值于 $F\otimes_{\Q_p}A$。

\begin{rem}{5.1.34}
提醒读者：虽然 $\Theta$ 的特征多项式总有系数属于 $A\otimes_{\Q_p}F$，这并不意味着可以用 $\Gamma$-作用把对 $(D_{\mathrm{Sen}}(D_A),\Theta)$（函子性地）下降到有限投射 $A\otimes_{\Q_p}F$-模连同其 $A\otimes_{\Q_p}F$-线性自同态。原因在于：连续的半线性 $\Gamma$-作用不是下降数据——只有光滑的半线性 $\Gamma$-作用才是。既然如此，人们或许会尝试用
\[
\Gamma\xrightarrow{\ \epsilon_{\mathrm{cyc}}\ }\Z_p^\times\xrightarrow{\ \log\ }\Z_p\longrightarrow(A\otimes_{\Q_p}F_\infty)^\times
\]
的逆来扭 $\Gamma$-作用，其中 $\epsilon_{\mathrm{cyc}}$ 是分圆特征标，最后一个映射由 $1\mapsto\exp(\Theta)$ 给出，以便得到光滑作用、从而是下降数据。然而这只有在 $\exp(\Theta)$ 收敛时才有意义，而一般情形并非如此。用 $p^n\Theta$（$n\gg0$）替换 $\Theta$ 总能保证收敛，但这样只能扭 $\Gamma$-作用在 $\Gamma_n\subset\Gamma$ 上的限制，并得到由 $F_\infty$ 到 $F_n$ 的下降数据。这一下降恰好给出上述讨论中 $A\otimes_{\Q_p}F_n$ 上的模 $D_{\mathrm{Sen},n}(D_A)$。
\label{thm:a4-5-1-34}
\end{rem}

若 $A$ 是 $L$-代数且 $|\Hom(F,L)|=[F:\Q_p]$，则 $F\otimes_{\Q_p}A\cong\prod_{\tau:F\hookrightarrow L}A_\tau$，其中 $A_\tau=A$。对每个嵌入 $\tau$，得到有限投射 $A\otimes_{\tau,F}F_\infty$-模
\[
D_{\mathrm{Sen},\tau}(D_A)=D_{\mathrm{Sen}}(D_A)\otimes_{A\otimes_{\Q_p}F}A_\tau
\]
的自同态 $\Theta_\tau$，其特征多项式 $\chi_{\Theta_\tau}$ 的系数属于 $A$。这个多项式的零点（即 $\Theta_\tau$ 的特征值）称为 $D$ 的 Hodge--Tate--Sen 权（Hodge--Tate--Sen weight），其标号为 $\tau$。

\ysrc{216}
若 $A=L$ 是包含 $\chi_{\Theta_\tau}$ 全部零点的域，则记 $\mathrm{wt}_\tau(D)\subseteq L$ 为这些零点之集。

记 $\mathfrak{gl}_n\sslash\mathrm{GL}_n$ 为 Lie 代数 $\mathfrak{gl}_n$ 在 $\mathrm{GL}_n$ 伴随作用下的 GIT 商。把 $\mathcal{R}_{F,A}$ 上的 $(\varphi,\Gamma)$-模 $D_A$ 送到 $\Theta$ 的特征多项式 $\chi_\Theta$ 的系数，定义了一个映射
\begin{equation}\tag{5.1.35}
\omega_d: \mathfrak{X}_d\longrightarrow\bigl(\mathrm{Res}_{F/\Q_p}(\mathfrak{gl}_n\sslash\mathrm{GL}_n)\bigr)_L=\mathrm{WT}_{d,L}.
\end{equation}
若 $|\Hom(F,L)|=[F:\Q_p]$，则当然有 $\mathrm{WT}_{d,L}=\prod_{\tau:F\to L}\mathrm{WT}_{\tau,d,L}$，且每个 $\mathrm{WT}_{\tau,d,L}=(\mathfrak{gl}_n\sslash\mathrm{GL}_n)_L$ 都是仿射空间。

本节最后，我们猜想：除经由行列式 $\mathrm{GL}_d\to\mathrm{GL}_1$ 诱导的映射拉回的函数之外，$\mathfrak{X}_d$ 上的整体函数都沿映射 $\omega_d$ 拉回得到。为给出精确表述，需要讨论行列式映射：对一个等变向量丛取最高外幂，定义了映射
\[
\det: \mathfrak{X}_d\longrightarrow\mathfrak{X}_1.
\]
由下文第 7.1 节的讨论（特别见 (7.1.3)），叠 $\mathfrak{X}_1$ 可写成商叠 $\mathcal{T}/\mathbb{G}_m$，其中 $\mathcal{T}$ 是 $F^\times$ 的连续特征标空间，配有刚性解析群 $\mathbb{G}_m$ 的平凡作用。由于 $\mathbb{G}_m$ 在 $\mathcal{T}$ 上的作用平凡，$\mathfrak{X}_1=\mathcal{T}\times */\mathbb{G}_m$ 有到 $\mathcal{T}$ 的典范投影；特别地，沿此映射的拉回诱导同构 $\Gamma(\mathfrak{X}_1,\mathcal{O}_{\mathfrak{X}_1})\cong\Gamma(\mathcal{T},\mathcal{O}_{\mathcal{T}})$。为简便，我们也常把 $\det$ 与这一投影的复合记作
\[
\det:\mathfrak{X}_d\longrightarrow\mathcal{T}.
\]
与上述构造类似，最高外幂定义了映射 $\mathrm{WT}_{d,L}\to\mathrm{WT}_{1,L}$，而 $\omega_1$ 诱导映射 $\mathcal{T}\to\mathrm{WT}_{1,L}$。

\begin{conj}{5.1.36}
映射 $\omega_d$ 与 $\det$ 诱导同构
{\small
\[
\begin{aligned}
\Gamma(\mathfrak{X}_d,\mathcal{O}_{\mathfrak{X}_d})
&\cong\Gamma(\mathrm{WT}_{d,L},\mathcal{O}_{\mathrm{WT}_{d,L}})\widehat{\otimes}_{\Gamma(\mathrm{WT}_{1,L},\mathcal{O}_{\mathrm{WT}_{1,L}})}\Gamma(\mathfrak{X}_1,\mathcal{O}_{\mathfrak{X}_1})\\
&\cong\Gamma(\mathrm{WT}_{d,L},\mathcal{O}_{\mathrm{WT}_{d,L}})\widehat{\otimes}_{\Gamma(\mathrm{WT}_{1,L},\mathcal{O}_{\mathrm{WT}_{1,L}})}\Gamma(\mathcal{T},\mathcal{O}_{\mathcal{T}}).
\end{aligned}
\]}
\label{thm:a4-5-1-36}
\end{conj}

刚性解析空间（或刚性解析空间上的叠）的连通分支与结构层整体截面的本原幂等元一一对应，因此猜想 5.1.36 给出了 $\mathfrak{X}_d$ 连通分支的描述。令
\[
\mathcal{W}=(\Spf\OO[[\OO_F^\times]])^{\mathrm{rig}}
\]
记 $\OO_F^\times$ 的连续特征标构成的刚性解析空间。于是 $F$ 的单值化元的选取诱导同构 $\mathcal{T}\cong\mathcal{W}\times\mathbb{G}_m$，我们由此得到 $\mathfrak{X}_d$ 连通分支的如下描述。

\begin{cor}{5.1.37}
假设猜想 5.1.36 成立。则 $\mathfrak{X}_d$ 的连通分支（经行列式）与 $\mathfrak{X}_1$ 的连通分支一一对应，从而与 $\mathcal{W}$ 的连通分支一一对应。
\label{thm:a4-5-1-37}
\end{cor}

注意 $\mathcal{W}$ 的连通分支又与特征标 $\mu(F)\to\OO^\times$ 一一对应，其中 $\mu(F)\subset F^\times$ 记单位根子群。还要指出，Böckle--Iyengar--Paškūnas [BIP23a] 对泛形变环的分支证明了类似结果。

\ysrc{217}
\subsection{de Rham 对象的模叠}

$(\varphi,\Gamma)$-模的叠 $\mathfrak{X}_d$ 有一些用 $p$-进 Hodge 理论定义的闭子叠，类似于闭子叠 $\mathcal{X}^{ss,\lambda,\tau}_d\subset\mathcal{X}_d$。在 Robba 环上的 $(\varphi,\Gamma)$-模情形，这些闭子叠可用（带滤过的）Weil--Deligne 表示来研究。在局部解析表示与叠 $\mathfrak{X}_d$ 上的层之间的猜想联系中，这些 de Rham 对象（给定权与惯性型）的闭子叠将在与光滑范畴 Langlands 猜想的比较中起作用。下面我们来定义这些「de Rham 轨迹」。

设 $F'$ 是 $F$ 的有限伽罗瓦扩张，伽罗瓦群为 $\Gal(F'/F)$；记 $F_0'$ 与 $F_0$ 分别为 $F'$ 与 $F$ 的极大非分歧子扩张。又设 $\sigma$ 记 Frobenius $x\mapsto x^p$ 到 $F_0'$ 的提升。我们考虑 $L$-概形范畴上的叠 $\mathrm{Mod}_{d,\varphi,N,F'/F}$，它把 $L$-代数 $A$ 送到秩 $d$ 的有限投射 $A\otimes_{\Q_p}F_0'$-模 $D$ 连同以下数据的群胚：
\begin{itemize}[leftmargin=2.2em,itemsep=0.2em]
\item 一个 $\mathrm{id}\otimes\sigma$-线性自同构 $\varphi_D: D\xrightarrow{\sim}D$；
\item 一个 $A\otimes_{\Q_p}F_0'$-线性自同态 $N: D\to D$，满足
\[
N\circ\varphi_D=p\varphi_D\circ\sigma^*N;
\]
\item $D$ 上的 $\Gal(F'/F)$-作用，与 $\varphi_D$、$N$ 交换，且满足
\[
g\bigl((a\otimes f)\cdot d\bigr)=\bigl(a\otimes g(f)\bigr)\cdot g(d)
\]
对所有 $a\in A$、$f\in F_0'$、$d\in D$、$g\in\Gal(F'/F)$ 成立。
\end{itemize}

我们还要回顾 $F$ 的 $d$ 维 Weil--Deligne 表示的叠 $\mathrm{WD}_{d,F}$：它是 $L$-概形范畴上的 Artin 叠，见 [DHKM25]、[Zhu25a, §3]。回顾：对给定的 $L$-概形 $S$，群胚 $\mathrm{WD}_{d,F}(S)$ 参数化 $S$ 上的向量丛 $\mathcal{E}$ 连同 Weil 群 $W_F$ 的光滑作用 $\rho$ 以及满足 $Ng=q^{\|g\|}gN$ 的自同态 $N:\mathcal{E}\to\mathcal{E}$。给定 $F'$，记 $\mathrm{WD}_{d,F'/F}\subset\mathrm{WD}_{d,F}$ 为 $\mathrm{WD}_{d,F}$ 的开闭子叠，由那些使 $\rho$ 在 $F'$ 的惯性子群 $I_{F'}$ 上的限制平凡的 Weil--Deligne 表示 $(\mathcal{E},\rho,N)$ 组成。

注意 Fontaine 的构造——把 $(\varphi,N,\Gal(F'/F))$-模对应到 Weil--Deligne 表示 [Fon94]（另见 [BS07, Prop. 4.1]）——表明这两个叠在足够大的域 $L$ 上同构：

\begin{lem}{5.2.1}
设 $L$ 足够大，使得 $[F_0':\Q_p]=|\Hom(F_0',L)|$，并固定嵌入 $F_0'\hookrightarrow L$。则存在叠的同构
\[
\mathrm{WD}_d: \mathrm{Mod}_{d,\varphi,N,F'/F}\xrightarrow{\ \sim\ }\mathrm{WD}_{d,F'/F}.
\]
\label{thm:a4-5-2-1}
\end{lem}

我们指出，该同构不是典范的，即它依赖于引理 5.2.1 中嵌入 $F_0'\hookrightarrow L$ 的选取。下文 §5.2.7 将详细写出这个同构，并给出所涉叠的显式描述（至少对 $F'=F$）。

为了把 $(\varphi,N,\Gal(F'/F))$-模 $D\in\mathrm{Mod}_{d,\varphi,N,F'/F}(A)$ 与等变向量丛联系起来，需要在 $D$ 上引入滤过。暂时记 $G=\mathrm{Res}_{F/\Q_p}\mathrm{GL}_d$、$G'=\mathrm{Res}_{F'/\Q_p}\mathrm{GL}_d$。此外记 $G_0'=\mathrm{Res}_{F_0'/\Q_p}\mathrm{GL}_d$。设 $T\subseteq G$ 记 $\mathrm{GL}_d$ 的对角环面的 Weil 限制，$B\subseteq G$ 记上三角矩阵的 Weil 限制。对支配余特征标 $\lambda\in X_*(T_{\bar L})_+$ 的选取，记 $L_\lambda$ 为 $\lambda$ 的反射域（reflex field），$P_\lambda\subseteq G_{L_\lambda}$ 为 $\lambda$ 定义的抛物子群（包含 $B_{L_\lambda}$）。

\ysrc{218}
旗簇 $G_{L_\lambda}/P_\lambda$ 于是参数化 $L_\lambda^d\otimes_{\Q_p}F$ 上型为 $\lambda$ 的滤过。更精确地说，利用标准等同
\[
X_*(T)_+=\prod_{\tau:F\hookrightarrow\bar L}\Z^d_+,
\]
我们写
\[
\lambda=(\lambda_{\tau,1}\ge\lambda_{\tau,2}\ge\cdots\ge\lambda_{\tau,d})_{\tau:F\hookrightarrow\bar L}.
\]
于是 $G_{L_\lambda}/P_\lambda$ 的 $\bar L$-值点由
\[
\bar L^d\otimes_{\Q_p}F=\prod_{\tau:F\hookrightarrow\bar L}\bar L^d_\tau
\]
（其中 $\bar L_\tau=\bar L$）上的一族递减、穷竭且分离的 $\Z$-滤过给出，使得 $(\lambda_{\tau,1}\ge\lambda_{\tau,2}\ge\cdots\ge\lambda_{\tau,d})$ 恰为 $\bar L^d_\tau$ 上滤过的权（计重数）。

利用典范映射 $G\to G'$，也可把 $\lambda$ 看作 $G'$ 中某个环面的余特征标（更精确地说，仍是 $\mathrm{GL}_d$ 中对角环面的 Weil 限制），从而得到类似的抛物子群 $P_\lambda'$ 与旗簇 $G'_{L_\lambda}/P_\lambda'$。

注意存在遗忘映射
\[
\mathrm{WD}_{d,F}\longrightarrow B\mathrm{GL}_d=*/\mathrm{GL}_d
\]
到秩 $d$ 向量丛的叠。类似地，存在典范映射
\[
\mathrm{Mod}_{d,\varphi,N,F'/F}\longrightarrow */G_0',
\]
它遗忘秩 $d$ 的投射 $A\otimes_{\Q_p}F_0'$-模上 $\varphi$、$N$ 与 $\Gal(F'/F)$ 的作用。

概形 $G/P_\lambda$ 配有 $\mathrm{GL}_d$ 的（对角）左平移作用，于是可定义纤维积
\[
\mathrm{Fil}_\lambda\mathrm{WD}_{d,F}=\mathrm{WD}_{d,F}\times_{*/\mathrm{GL}_d}\mathrm{GL}_d\backslash G_{L_\lambda}/P_\lambda,
\]
它称为（Hodge--Tate 权为 $\lambda$ 的）带滤过 Weil--Deligne 表示的叠。其 $S$-值点参数化 $S$ 上向量丛 $\mathcal{E}$ 上的 Weil--Deligne 表示，连同 $F\otimes_{\Q_p}\mathcal{E}$ 上型为 $\lambda$ 的滤过。

类似地，概形 $G'/P_\lambda'$ 配有 $G_0'$ 的（对角）左平移作用，我们可以考虑闭子叠
{\small
\[
\mathrm{Fil}_\lambda\mathrm{Mod}_{d,\varphi,N,F'/F}\subset\mathrm{Mod}_{d,\varphi,N,F'/F}\times_{*/G_0'}G_0'\backslash G'_{L_\lambda}/P_\lambda',
\]}
它（对给定对象 $D\in\mathrm{Mod}_{d,\varphi,N,F'/F}$）参数化 $D\otimes_{F_0'}F'$ 上型为 $\lambda$ 且在 $\Gal(F'/F)$ 作用下稳定的那些滤过。

\begin{rem}{5.2.2}
我们指出，$\Gal(F'/F)$-稳定滤过只可能在 $G'$ 的（我们的极大环面的）余特征标 $\lambda$ 经由映射 $G\to G'$ 来自 $G$ 的（我们的极大环面的）余特征标时才存在。
\label{thm:a4-5-2-2}
\end{rem}

下述引理推广了引理 5.2.1。

\begin{lem}{5.2.3}
设 $L$ 足够大，使得 $[F':\Q_p]=|\Hom(F',L)|$，并固定嵌入 $F'\hookrightarrow L$。则存在典范同构
\[
\mathrm{Fil}_\lambda\mathrm{Mod}_{d,\varphi,N,F'/F}\xrightarrow{\ \sim\ }\mathrm{Fil}_\lambda\mathrm{WD}_{d,F'/F},
\]
它位于引理 5.2.1 的同构 $\mathrm{WD}_d$ 之上。
\label{thm:a4-5-2-3}
\end{lem}

\ysrc{219}
叠 $\mathrm{Fil}_\lambda\mathrm{Mod}_{d,\varphi,N,F'/F}$ 在 de Rham $(\varphi,\Gamma)$-模的刻画中起重要作用；经由上述同构，我们于是能把带滤过 Weil--Deligne 表示与 $(\varphi,\Gamma)$-模联系起来——至少在基域足够大时。更精确地，有下述定理。

\begin{thm}{5.2.4}
对固定的 $\lambda\in X_*(T_{\bar L})_+$，存在典范闭嵌入
\[
(\mathrm{Fil}_\lambda\mathrm{Mod}_{d,\varphi,N,F'/F})^{\mathrm{an}}\longrightarrow\mathfrak{X}_{d,L_\lambda}.
\]
此外，该态射的像恰为 Hodge--Tate--Sen 权为 $\lambda$、且在基变换到 $F'$ 后成为半稳定的 de Rham $(\varphi,\Gamma)$-模构成的闭子叠。
\label{thm:a4-5-2-4}
\end{thm}

在刚性解析叠 $(\mathrm{Fil}_\lambda\mathrm{Mod}_{d,\varphi,N,F'/F})^{\mathrm{an}}$ 上构造秩 $d$ 的 $(\varphi,\Gamma)$-模族，是把 Berger [Ber08b] 的工作推广到所涉对象的族。

为理解定理 5.2.4 的「此外」部分，需要定义 $\mathcal{R}_{F,A}$ 上（潜在）半稳定或 de Rham $(\varphi,\Gamma)$-模族的概念。给定 $\mathcal{R}_{F,A}$ 上的 $(\varphi,\Gamma)$-模 $D_A$，可考虑基变换
\[
D_{\mathrm{dif}}(D_A)=D_A\otimes_{\mathcal{R}_{F,A}}F_\infty(\!(t)\!),
\]
它带有位于 $t\frac{d}{dt}$ 之上的导子
\[
\partial: D_{\mathrm{dif}}(D_A)\longrightarrow D_{\mathrm{dif}}(D_A),
\]
该导子由 $\Gamma$-作用的导数定义。于是称 $D_A$ 为 de Rham 的，若 $D_{\mathrm{dif}}(D_A)^{\partial=0}$ 是秩 $d$ 的投射 $A\otimes_{\Q_p}F_\infty$-模，且典范映射
\[
D_{\mathrm{dif}}(D_A)^{\partial=0}\otimes_{F_\infty}F_\infty(\!(t)\!)\longrightarrow D_{\mathrm{dif}}(D_A)
\]
是同构。为证明定理 5.2.4 的「此外」部分，我们紧循 [BC08, Thms. 5.3.1, 5.3.2] 的证明，来说明 $\mathfrak{X}_{d,L_\lambda}$ 中存在由 Hodge--Tate 权为 $\lambda$ 的 de Rham 对象组成的 Zariski 闭子叠。由 $p$-进单值性定理（$p$-adic monodromy theorem），经典点（即对 $L_\lambda$ 的有限扩张 $L'$ 的 $L'$-值点）与 $\bigcup_{F'}(\mathrm{Fil}_\lambda\mathrm{Mod}_{d,\varphi,N,F'/F})^{\mathrm{an}}$ 的经典点重合，这说明这两个叠必须重合。

\begin{rem}{5.2.5}
回顾：解析设定中的 $(\varphi,\Gamma)$-模不要求是 étale 的。因此 de Rham $(\varphi,\Gamma)$-模的叠可以用带滤过 $(\varphi,N)$-模（带下降数据）来描述，正如定理 5.2.4 所做的那样，而不必对滤过附加「弱容许性（weak admissibility）」条件。不过，可以过渡到 $(\mathrm{Fil}_\lambda\mathrm{Mod}_{d,\varphi,N,F'/F})^{\mathrm{an}}$ 的一个开子叠，在其上联系于带滤过 $(\varphi,N,\Gal_F)$-模的 $(\varphi,\Gamma)$-模是 étale 的；$\Gal_F$-作用平凡情形的这一构造见 [HH20] 与 [Hel13]。这个 étale 轨迹随后与叠 $\mathcal{X}^{ss,\lambda,\tau}_d$ 在映射 (5.1.31) 下的一般纤维之像（对适当的 $\tau$ 取并）一致。
\label{thm:a4-5-2-5}
\end{rem}

\begin{rem}{5.2.6}
我们指出，定理中固定余特征标 $\lambda$ 是重要的。若允许 $\lambda$ 任意变动，则相应子叠 $(\mathrm{Fil}_\lambda\mathrm{Mod}_{d,\varphi,N,F'/F})^{\mathrm{an}}$ 之并在 $\mathfrak{X}_d$ 中不闭，而预期是 Zariski 稠密的。这应视为晶体表示在局部形变环中稠密这一事实的类比，见 [Che13, Nak14, BIP23b]。
\label{thm:a4-5-2-6}
\end{rem}

\ysrc{220}
\subsubsection*{5.2.7\quad 嵌入 $F_0'\hookrightarrow L$ 的作用}

在引理 5.2.1 中，我们在包含 $F_0'$ 的域 $L$ 上，借助嵌入 $\sigma_0: F_0'\hookrightarrow L$ 的选取，把叠 $\mathrm{Mod}_{d,\varphi,N,F'/F}$ 与 $\mathrm{WD}_{d,F'/F}$ 等同起来。下面进一步说明这个嵌入的选取。为简单起见（也为了避免记号过载），只讨论 $F'=F$ 的情形。下文所用叠的显式描述参见 [BG19, §2.6] 与 [HH20, §3]。

设 $L$ 包含 $F_0'=F_0$（等式成立是因为已假设 $F'=F$）；则对 $L$-代数 $A$ 有分解
\[
A\otimes_{\Q_p}F_0=A\otimes_L(L\otimes_{\Q_p}F_0)=A\otimes_L\prod_{\varsigma:F_0\hookrightarrow L}L_\varsigma=\prod_{\varsigma:F_0\hookrightarrow L}A_\varsigma,
\]
其中对一切 $\varsigma$ 有 $L_\varsigma=L$、$A_\varsigma=A$。$\Gal(F_0/\Q_p)$ 的作用当然只是置换乘积中的因子。本节余下部分固定一个嵌入 $\varsigma_0$，并把 $A\otimes_{\Q_p}F_0=\prod_{i=1}^fA_i$ 重写，其中 $A_i=A_{\varsigma_i^0}=A$。由此得到同构
\begin{align}
(\mathrm{Res}_{F_0/\Q_p}\mathrm{GL}_d)_L&=\prod_{i=1}^f\mathrm{GL}_d,\tag{5.2.8}\\
(\mathrm{Res}_{F_0/\Q_p}\mathrm{Mat}_{d\times d})_L&=\prod_{i=1}^f\mathrm{Mat}_{d\times d}.\notag
\end{align}

注意 $\mathrm{WD}_{d,F/F}=X_d/\mathrm{GL}_d$，其中
\[
X_d=\{(\varphi,N)\in\mathrm{GL}_d\times\mathrm{Mat}_{d\times d}\mid N\varphi=q\varphi N\},
\]
且 $\mathrm{GL}_d$ 通过共轭作用在 $\mathrm{GL}_d$ 与 $\mathrm{Mat}_{d\times d}$ 上。另一方面
\[
\mathrm{Mod}_{d,\varphi,N,F/F}=Y_d/\mathrm{Res}_{F_0/\Q_p}\mathrm{GL}_d,
\]
其中
\[
Y_d=\{(\varphi,N)\in\mathrm{Res}_{F_0/\Q_p}\mathrm{GL}_d\times\mathrm{Res}_{F_0/\Q_p}\mathrm{Mat}_{d\times d}\mid N\varphi=p\varphi\sigma^*N\},
\]
而 $\sigma:\mathrm{Res}_{F_0/\Q_p}\mathrm{Mat}_{d\times d}\to\mathrm{Res}_{F_0/\Q_p}\mathrm{Mat}_{d\times d}$ 是由 $F_0$ 上的 $\sigma$ 诱导的映射；$\mathrm{Res}_{F_0/\Q_p}\mathrm{GL}_d$ 通过共轭作用在 $\mathrm{Res}_{F_0/\Q_p}\mathrm{Mat}_{d\times d}$ 上，并通过 $\sigma$-共轭作用在 $\mathrm{Res}_{F_0/\Q_p}\mathrm{GL}_d$ 上，即 $g\cdot(\varphi,N)=(g^{-1}\varphi\sigma(g),g^{-1}Ng)$。

基变换到 $L$ 后，用 (5.2.8) 把 $Y_d$ 的点写成多元组 $((\varphi_1,\dots,\varphi_f),(N_1,\dots,N_f))$，其中 $\varphi_i\in\mathrm{GL}_d$、$N_i\in\mathrm{Mat}_{d\times d}$，且
\[
N_{i+1}\varphi_i=p\varphi_iN_i,\qquad i=1,\dots,f,
\]
这里约定 $N_{f+1}=N_1$。$(\mathrm{Res}_{F_0/\Q_p}\mathrm{GL}_d)_L\cong\prod_{i=1}^f\mathrm{GL}_d$ 的作用于是由
\[
\begin{aligned}
&(g_1,\dots,g_f)\cdot\bigl((\varphi_1,\dots,\varphi_f),(N_1,\dots,N_f)\bigr)\\
&\qquad=\bigl((g_1^{-1}\varphi_1g_f,\ g_2^{-1}\varphi_2g_1,\ \dots,\ g_f^{-1}\varphi_fg_{f-1}),\ (g_1^{-1}N_1g_1,\ \dots,\ g_f^{-1}N_fg_f)\bigr)
\end{aligned}
\]
给出。

利用这一显式描述，可写出典范态射 $X_d\to Y_d$（它依赖于同构 (5.2.8)，从而依赖于 $\varsigma_0$ 的选取）：
\[
(\varphi,N)\longmapsto\bigl((\varphi,1,\dots,1),(N,p\varphi N\varphi^{-1},p^2\varphi N\varphi^{-1},\dots,p^{(f-1)}\varphi N\varphi^{-1})\bigr),
\]
它对群概形态射 $\mathrm{GL}_d\to\mathrm{Res}_{F_0/\Q_p}\mathrm{GL}_d$（$g\mapsto(g,g,\dots,g)$）等变，并诱导叠的同构
\[
f_{\varsigma_0}: \mathrm{WD}_{d,F/F}=X_d/\mathrm{GL}_d\xrightarrow{\ \sim\ }\mathrm{Mod}_{d,\varphi,N,F/F}=Y_d/\mathrm{Res}_{F_0/\Q_p}\mathrm{GL}_d.
\]

\ysrc{221}
这个同构的逆由映射 $Y_d\to X_d$ 诱导，后者由
\[
\bigl((\varphi_1,\dots,\varphi_f),(N_1,\dots,N_f)\bigr)\longmapsto(\varphi_f\cdots\varphi_2\varphi_1,N_1)
\]
给出。

这些公式表明，同构 $f_{\varsigma_0}$ 依赖于 $\varsigma_0$ 的选取：对另一个选取 $\varsigma_0'$，诱导的自同构 $f_{\varsigma_0'}\circ f_{\varsigma_0}^{-1}$ 不是恒等。此外容易看出，$f_{\varsigma_0}$ 不能下降到不含 $F_0$ 的子域 $L'\subset L$：若 $\tau:L\to L$ 是 $L$ 的域自同构，则它与 $f_{\varsigma_0}$ 交换当且仅当它在 $F_0$ 上的限制是恒等。

另一方面注意，映射
\[
\mathrm{WD}_{d,F/F}\longrightarrow\mathrm{GL}_d\sslash\mathrm{GL}_d\cong\mathbb{A}^{d-1}\times\mathbb{G}_m
\]
把 $(\varphi,N)$ 送到其特征多项式的系数，它在整体截面上诱导同构，例如见 [Hel23, Remark 3.4]。由此推出，$f_{\sigma_0}$\tn{底本此处作 $f_{\sigma_0}$，疑为 $f_{\varsigma_0}$ 之误。} 诱导结构层的（不变）整体截面上的同构
\begin{equation}\tag{5.2.9}
\Gamma(\mathrm{Mod}_{d,\varphi,N,F/F},\mathcal{O}_{\mathrm{Mod}_{d,\varphi,N,F/F}})\longrightarrow\Gamma(\mathrm{WD}_{d,F/F},\mathcal{O}_{\mathrm{WD}_{d,F/F}}),
\end{equation}
它与 $\varsigma_0$ 的选取无关；而且该同构可下降到任意子域 $L'\subseteq L$。

下面在 $L=F_0$、$L'=\Q_p$ 的情形把这一下降显式化。此时 $\Gal(F_0/\Q_p)$ 的生成元 $\sigma$ 在 $Y_d$ 的 $\varphi$-部分上的作用 $\sigma_{Y_d}$ 由
\[
(\varphi_1,\dots,\varphi_f)\longmapsto(\sigma^*\varphi_f,\sigma^*\varphi_1,\dots,\sigma^*\varphi_{f-1})
\]
给出，而 $\sigma$ 在 $X_d$ 的 $\varphi$-部分上的作用 $\sigma_{X_d}$ 显然由 $\varphi\mapsto\sigma^*\varphi$ 给出。由此推出（只写出公式的 $\varphi$-部分）
\[
(f_{\varsigma_0}\circ\sigma_{Y_d})(\varphi_1,\dots,\varphi_f)=(\sigma^*\varphi_f)^{-1}\bigl(\sigma_{X_d}\circ f_{\varsigma_0}\bigl((\varphi_1,\dots,\varphi_f)\bigr)\bigr)(\sigma^*\varphi_f),
\]
从而 $f_{\varsigma_0}$ 不能下降到 $\Q_p$。另一方面，$f_{\varsigma_0}\circ\sigma_{Y_d}$ 与 $\sigma_{X_d}\circ f_{\varsigma_0}$ 的 $\varphi$-部分的特征多项式一致，因此整体截面的同构 (5.2.9) 可下降到 $\Q_p$。

事实上，这一观察对 $\mathrm{WD}_{d,F'/F}$ 有如下推广。

\begin{prop}{5.2.10}
设 $L$ 是 $\Q_p$ 的任意有限扩张，$F'/F$ 是有限扩张。对 $L$ 的包含 $F_0'$ 的域扩张 $M$ 以及嵌入 $\varsigma_0: F_0'\hookrightarrow M$ 的任意选取，由 $f_{\varsigma_0}$ 诱导的同构
\[
\Gamma\bigl((\mathrm{Mod}_{d,\varphi,N,F'/F})_M,\mathcal{O}_{(\mathrm{Mod}_{d,\varphi,N,F'/F})_M}\bigr)\longrightarrow\Gamma\bigl((\mathrm{WD}_{d,F'/F})_M,\mathcal{O}_{(\mathrm{WD}_{d,F'/F})_M}\bigr)
\]
与 $\varsigma_0$ 的选取无关，并可下降为典范同构
\[
\Gamma(\mathrm{Mod}_{d,\varphi,N,F'/F},\mathcal{O}_{\mathrm{Mod}_{d,\varphi,N,F'/F}})\longrightarrow\Gamma(\mathrm{WD}_{d,F'/F},\mathcal{O}_{\mathrm{WD}_{d,F'/F}}).
\]
\label{thm:a4-5-2-10}
\end{prop}

\subsubsection*{5.2.11\quad 不变函数与 Bernstein 中心}

与 §4.1 中按给定的惯性型 $\tau$ 定义 $\mathcal{X}^{ss,\lambda,\tau}_d$ 的方式类似，我们将定义 Hodge--Tate 权为 $\lambda$、惯性型 $\tau$ 固定的 de Rham 对象叠。回顾 $\mathrm{GL}_d(F)$ 的局部 Langlands 对应给出双射 $\tau\mapsto\Omega_\tau$，把 $I_F$ 的（定义在 $L$ 上的）$d$ 维惯性型 $\tau$ 与范畴 $\Rep^{\mathrm{sm}}\mathrm{GL}_d(F)$ 的（定义在 $L$ 上的）Bernstein 块 $\Omega$ 对应起来（我们把表示看作 $L$ 的代数闭包 $\bar L$ 上的向量空间上的表示）。更精确地说，若 $(r,N)$ 是联系于 $\mathrm{GL}_d(F)$ 的不可约光滑表示 $\pi$ 的 Weil--Deligne 表示，则 $r|_{I_F}\cong\tau$ 当且仅当 $\pi$ 落在 Bernstein 块 $\Omega_\tau$ 中。（这是惯性局部 Langlands 对应的一种形式；在 $\mathrm{GL}_2(\Q_p)$ 的情形，这一点已在 §3.2.1 中指出。此双射 $\tau\leftrightarrow\Omega_\tau$ 与用 $K$-型表述的惯性局部 Langlands 之间的关系，将在下文 §6.1.9 中回顾。）

\ysrc{222}
此外注意，在 $L$ 的代数闭包 $\bar L$ 上存在双射 $\tau\mapsto\mathrm{WD}_{d,F,\tau}$，把（定义在 $\bar L$ 上的）$d$ 维惯性型 $\tau$ 与 Weil--Deligne 表示的叠的连通分支对应起来，使得 Weil--Deligne 表示 $(r,N)$ 落在 $\mathrm{WD}_{d,F,\tau}$ 中当且仅当 $r|_{I_F}\cong\tau$，见 [BZCHN24, §5.2]。由于我们只限于在 $I_{F'}$ 上变成平凡的惯性型，可以取 $L$ 足够大，使得 $\bar L$ 上所有惯性型 $\tau$（满足 $\tau|_{I_{F'}}$ 平凡）都定义在 $L$ 上，且所有几何连通分支 $\mathrm{WD}_{d,F,\tau}$ 都定义在 $L$ 上；我们（稍滥用地）仍把这个 $L$-概形上的叠记作 $\mathrm{WD}_{d,F,\tau}$。

下述命题（它是 [BZCHN24, Theorem 5.13] 的推论，整系数版本见 [HM18]）把 $\Omega_\tau$ 的 Bernstein 中心 $Z_{\Omega_\tau}$（把它看作 $L$-向量空间上光滑表示范畴的一个块，特别地 $Z_{\Omega_\tau}$ 是 $L$-代数）与 $\mathrm{WD}_{d,F,\tau}$ 上的不变函数联系起来。

\begin{prop}{5.2.12}
设 $\tau$ 是定义在 $L$ 上的惯性型。则存在自然同构
\[
Z_{\Omega_\tau}\xrightarrow{\ \sim\ }\Gamma(\mathrm{WD}_{d,F,\tau},\mathcal{O}_{\mathrm{WD}_{d,F,\tau}}).
\]
\label{thm:a4-5-2-12}
\end{prop}

\begin{rem}{5.2.13}
注意命题 5.2.12 的同构有两种可能的正规化，对应于局部 Langlands 对应的两种可能的正规化。在通常的（酉）正规化中，光滑表示 $\pi$ 的中心特征与相应 Weil--Deligne 表示 $(\rho,N)$ 的行列式相匹配。有时更方便的是选取使 $(\rho,N)$ 对应 $\pi\otimes|\det|^{-\frac{d-1}{2}}$ 的正规化。这种正规化更适合局部–整体相容性，其优点是不依赖于 $q^{1/2}$ 的选取（因此在非代数闭域上考虑时性质更好，比较 [BS07]）。本讲义此后一律采用非酉正规化。
\label{thm:a4-5-2-13}
\end{rem}

注意，给定 $\tau$，引理 5.2.1 的同构使我们能定义 $\mathrm{Mod}_{d,\varphi,N,F'/F}$ 中对应于 $\mathrm{WD}_{d,F,\tau}$ 的连通分支 $\mathrm{Mod}_{d,\varphi,N,\tau}$。该连通分支不依赖于引理 5.2.1 中嵌入 $\varsigma_0$ 的选取。我们进一步用它定义 $\mathrm{Fil}_\lambda\mathrm{Mod}_{d,\varphi,N,F'/F}$ 的显然连通分支 $\mathrm{Fil}_\lambda\mathrm{Mod}_{d,\varphi,N,\tau}$。

对固定的惯性型 $\tau$，可考虑 §4 中定义的形变叠 $\mathcal{X}^{ss,\lambda,\tau}_d$ 的刚性解析一般纤维（generic fiber）$(\mathcal{X}^{ss,\lambda,\tau}_d)^{\mathrm{rig}}_\eta$。由 $\mathcal{X}^{ss,\lambda,\tau}_d$ 的构造本身，该一般纤维带有映射
\begin{equation}\tag{5.2.14}
\pi^{ss,\lambda,\tau}_d: (\mathcal{X}^{ss,\lambda,\tau}_d)^{\mathrm{rig}}_\eta\longrightarrow(\mathrm{Fil}_\lambda\mathrm{Mod}_{d,\varphi,N,\tau})^{\mathrm{an}}\subset\mathfrak{X}_d,
\end{equation}
它是 Pappas--Rapoport 的周期映射（period map）[PR09, (5.37)] 的推广，后者对应 $\tau=1$、$\lambda$ 极小的情形。此外它使下图交换：
\[
\begin{tikzcd}[column sep=large,row sep=large]
(\mathcal{X}^{ss,\lambda,\tau}_d)^{\mathrm{rig}}_\eta \arrow[r] \arrow[d,"\pi^{ss,\lambda,\tau}_d"']
& \mathfrak{X}^{\mathrm{rig}}_{d,\eta} \arrow[d,"\pi_d"]\\
(\mathrm{Fil}_\lambda\mathrm{Mod}_{d,\varphi,N,\tau})^{\mathrm{an}} \arrow[r]
& \mathfrak{X}_d
\end{tikzcd}
\]

\ysrc{223}
其中 $\pi_d$ 是 (5.1.31) 中的映射。特别地，由命题 5.2.10 与 5.2.12 合并得到的同构
\begin{equation}\tag{5.2.15}
Z_{\Omega_\tau}\xrightarrow{\ \sim\ }\Gamma(\mathrm{Mod}_{d,\varphi,N,\tau},\mathcal{O}_{\mathrm{Mod}_{d,\varphi,N,\tau}})
\end{equation}
诱导映射
\begin{equation}\tag{5.2.16}
Z_{\Omega_\tau}\longrightarrow\Gamma\bigl((\mathcal{X}^{ss,\lambda,\tau}_d)^{\mathrm{rig}}_\eta,\mathcal{O}_{(\mathcal{X}^{ss,\lambda,\tau}_d)^{\mathrm{rig}}_\eta}\bigr),
\end{equation}
它在使 $\tau$ 有定义的任意选定基域 $L$ 上都有定义（不必包含那个使我们能把 $(\varphi,N,\Gal(F'/F))$-模的叠与 Weil--Deligne 表示的叠的某个子叠等同的域 $F_0'$）。

\begin{rem}{5.2.17}
与到潜在半稳定形变环一般纤维上的函数环的映射复合之后，该映射与 [CEG$^+$16, Theorem 4.1, Prop. 4.2] 中的态射一致。事实上 [CEG$^+$16] 的结果更强：它断言该映射取值于有界函数环
\[
\Gamma(\mathcal{X}^{ss,\lambda,\tau}_d,\mathcal{O}_{\mathcal{X}^{ss,\lambda,\tau}_d})[\tfrac{1}{p}]\subset\Gamma\bigl((\mathcal{X}^{ss,\lambda,\tau}_d)^{\mathrm{rig}}_\eta,\mathcal{O}_{(\mathcal{X}^{ss,\lambda,\tau}_d)^{\mathrm{rig}}_\eta}\bigr).
\]
\label{thm:a4-5-2-17}
\end{rem}

\subsection{Drinfeld 型紧化}

我们讨论叠 $\mathfrak{X}_d$ 的若干变体，它们在一族秩 $d$ 的等变向量丛上参数化旗。这些叠自然地联系于抛物子群 $P\subseteq G=\mathrm{GL}_d$ 的选取。给定 Levi 商为 $M\cong\mathrm{GL}_{d_1}\times\cdots\times\mathrm{GL}_{d_r}$ 的抛物子群 $P$，定义 $\mathfrak{X}_P$ 为 Fargues--Fontaine 曲线 $X_{\bar F}$ 上 $\Gal_F$-等变 $P$-丛的叠，它带有两个投影
\[
\begin{tikzcd}[column sep=huge,row sep=large]
& \mathfrak{X}_P \arrow[dl,"\beta_P"'] \arrow[dr,"\alpha_P"] &\\
\mathfrak{X}_d & & \mathfrak{X}_M
\end{tikzcd}
\]
分别由映射 $P\hookrightarrow\mathrm{GL}_d$ 与 $P\twoheadrightarrow M$ 诱导，其中 $\mathfrak{X}_M\cong\mathfrak{X}_{d_1}\times\cdots\times\mathfrak{X}_{d_r}$ 记 $X_{\bar F}$ 上 $\Gal_F$-等变 $M$-丛的叠。下文方便起见记 $\mathfrak{X}_G=\mathfrak{X}_d$。

在 $P=B$ 的特殊情形，态射 $\mathfrak{X}_B\to\mathfrak{X}_d$ 的像中的 $(\varphi,\Gamma)$-模，按定义恰为三角化（trianguline）的 $(\varphi,\Gamma)$-模（Colmez 意义下；关于三角化表示的定义与讨论及其在 $p$-进 Langlands 对应中的作用，见 [Col08, Col14]）。三角化 $(\varphi,\Gamma)$-模本身是极自然的对象（三角化无非就是一个 $B$-结构），除此之外，其定义的动机还在于：联系于有限斜率的超收敛 $p$-进模形式（即本征系统由 Coleman--Mazur 特征曲线参数化的那些 $p$-进模形式）的伽罗瓦表示在 $p$ 处是三角化的，这是 Kisin [Kis03] 发现的（当时「三角化」一词尚未出现，且用的是等价描述）。这一结果后来被推广到更一般的特征簇（特别是酉群的特征簇，例如见 [BC09]）。我们将在第 9.6 节回到这一点，讨论特征簇、$p$-进自守形式与 $\mathfrak{X}_d$ 上的层之间的关系，后者用 $\mathfrak{X}_B$ 及其紧化来定义。

\begin{rem}{5.3.1}
当然存在各种变体，把 $\mathrm{GL}_d$ 换成任意约化群；但为陈述简便，这里只限于 $\mathrm{GL}_d$ 的情形。我们定义这些叠的方式使得一般情形下定义该取何形状一目了然。
\label{thm:a4-5-3-1}
\end{rem}

\ysrc{224}
追随 Drinfeld，Braverman 与 Gaitsgory [BG02] 对代数曲线上 $P$-丛的叠定义了两类所谓的 Drinfeld 紧化；下文我们将在 $\Gal_F$-等变向量丛的框架中、等价地在 $(\varphi,\Gamma)$-模的框架中，定义它们的类似物 $\widetilde{\mathfrak{X}}_P$ 与 $\overline{\mathfrak{X}}_P$。这两种紧化各有优劣：紧化 $\widetilde{\mathfrak{X}}_P$ 的优点是对抛物子群的包含 $P'\subseteq P$ 具有函子性，缺点是映射 $\alpha_P: \mathfrak{X}_P\to\mathfrak{X}_M$ 不能延拓到 $\widetilde{\mathfrak{X}}_P$。反之，紧化 $\overline{\mathfrak{X}}_P$ 的优点是 $\alpha_P$ 可延拓为 $\overline{\alpha}_P: \overline{\mathfrak{X}}_P\to\mathfrak{X}_M$，但 $P\mapsto\overline{\mathfrak{X}}_P$ 对抛物子群的包含 $P'\subseteq P$ 不具函子性。

在给出正式定义之前，先简要说明 $P=B$ 情形的动机。$(\varphi,\Gamma)$-模 $D_A$ 上的 $B$-结构是一个完全的 $(\varphi,\Gamma)$-稳定旗
\[
0=\mathrm{Fil}_0\subseteq\mathrm{Fil}_1\subseteq\cdots\subseteq\mathrm{Fil}_d=D_A,
\]
使得 $\mathrm{Fil}_i/\mathrm{Fil}_{i-1}$ 是秩 $1$ 的投射 $\mathcal{R}_{F,A}$-模。用 Plücker 坐标来说，这样的滤过数据等同于 $(\varphi,\Gamma)$-稳定的秩 $1$ 子对象
\[
\mathcal{L}_i=\bigwedge^i\mathrm{Fil}_i\subseteq\bigwedge^i D_A
\]
的数据，使得商 $(\bigwedge^iD_A)/\mathcal{L}_i$ 是投射 $\mathcal{R}_{F,A}$-模。这些直线 $\mathcal{L}_i$ 满足若干关系（称为 Plücker 关系），我们不把它们显式写出，而引用下文的更抽象描述——后者能更统一地描述这些关系（代价是要考虑 $D_A$ 在 $\mathrm{GL}_d$ 的所有代数表示上的取值，而不只是外幂 $\bigwedge^iD_A$，另比较下文注 5.3.5）。

下述例子说明：确实存在刚性空间 $\operatorname{Sp} A$ 上的族 $D_A$，它在某个 Zariski 开稠密子空间上有 $B$-结构，但该 $B$-结构不能延拓到整个 $\operatorname{Sp} A$。因此我们希望以某种方式紧化这一状况。想法是仍然参数化满足 Plücker 关系的直线 $\mathcal{L}_i\subseteq\bigwedge^iD_A$，但放弃 $(\bigwedge^iD_A)/\mathcal{L}_i$ 在 $\mathcal{R}_{F,A}$ 上投射这一条件。这显然解决了下面例子中的问题。

\begin{exam}{5.3.2}
设 $F=\Q_p$，考虑射影直线 $\PP^1$ 上下述带滤过 $\varphi$-模族。令
\[
M=\mathcal{O}_{\PP^1}e_1\oplus\mathcal{O}_{\PP^1}e_2,\qquad\varphi_M=\mathrm{diag}(\varphi_1,\varphi_2),
\]
其中 $\varphi_1\neq\varphi_2\in\Q_p^\times$ 固定。又考虑 $M$ 上的 $\Z$-滤过 $\mathrm{Fil}^\bullet$：
\[
\mathrm{Fil}^iM=\begin{cases}M & i\le-1\\ \mathscr{L} & i=0\\ 0 & i\ge1,\end{cases}
\]
其中 $\mathscr{L}\subset M$ 是 $\PP^1$ 上的泛直线。考虑子模 $M'=\mathcal{O}_{\PP^1}e_1\subset M$。则不存在 $M'$ 上的滤过 $\mathrm{Fil}^iM'$（含义是 $\mathrm{Fil}^iM'\subset M'$ 在 $\PP^1$ 上局部是直和项），使得对一切 $x\in\PP^1$，子对象 $\mathrm{Fil}^iM'\otimes k(x)\subset M'\otimes k(x)$ 都是 $\mathrm{Fil}^iM\otimes k(x)$ 与 $M'\otimes k(x)$ 的交。若坚持 $\mathrm{Fil}^iM'\subset M'$（局部）是直和项，则只能在 $\mathbb{A}^1\subset\PP^1$ 上做到这一点，即令
\[
\mathrm{Fil}^iM'=\begin{cases}M' & i\le-1\\ 0 & i\ge0.\end{cases}
\]
但在 $x=\infty$ 处，特殊化 $\mathrm{Fil}^iM'\otimes k(x)$ 是交 $\mathrm{Fil}^iM\otimes k(x)\cap M'\otimes k(x)$ 的真子集。
\label{thm:a4-5-3-2}
\end{exam}

\ysrc{225}
利用由带滤过 $\varphi$-模构造 $(\varphi,\Gamma)$-模的方法（见定理 5.2.4 及其后的讨论），上述情形翻译为如下状况。令
\[
\mathcal{R}^{\oplus2}_{F,\PP^1}\subset D\subset t^{-1}\mathcal{R}^{\oplus2}_{F,\PP^1}
\]
为 $\PP^1$ 上联系于 $M$ 的 $(\varphi,\Gamma)$-模，$\mathcal{L}$ 为子带滤过 $\varphi$-模 $M'\subset M$ 所定义的 $(\varphi,\Gamma)$-稳定直线。则 $\mathcal{L}\subset D$ 是 $D|_{\mathbb{A}^1}$ 上的 $B$-结构，但特殊化到 $\infty\in\PP^1$ 之后，该包含的余核有 $t$-挠。

为一般地定义这些紧化，引入下述记号，参 [BG02]。固定 Borel 子群 $B\subseteq G=\mathrm{GL}_d$，记 $T=B/U$，其中 $U$ 是 $B$ 的幂幺根（unipotent radical）。按惯例，$\langle\cdot,\cdot\rangle$ 记 $X^*(T)$ 与 $X_*(T)$ 之间的典范配对。给定抛物子群 $B\subseteq P\subseteq G$，其幂幺根为 $U_P$、Levi 商为 $M=P/U_P$，记 $X^*(T)^+_M\supset X^*(T)^+_G$ 分别为 $M$ 与 $G$ 的支配权（dominant weight）半群。$G$ 的单余根之集 $\{\alpha_i\in X_*(T),i\in\mathcal{I}\}$ 以 $G$ 的 Dynkin 图的顶点集 $\mathcal{I}=\mathcal{I}_G=\{1,\dots,d-1\}$ 为下标。对上述 $M$，$M$ 的 Dynkin 图给出子集 $\mathcal{I}_M\subseteq\mathcal{I}$，并令
\[
X^*(T)_{G,P}=\{\check\lambda\in X^*(T)\mid\langle\alpha_i,\check\lambda\rangle=0\ \text{对}\ i\in\mathcal{I}_M\}.
\]
此外记 $S_M=M/[M,M]$ 为 $M$ 的极大环面商。

给定 $X_{\bar F}$ 上（一族）$\Gal_F$-等变向量丛 $\mathcal{E}$（或 $\mathcal{R}_F$ 上的一族 $(\varphi,\Gamma)$-模 $D$）以及 $G$ 的代数表示 $V$，记 $V(\mathcal{E})$（相应地 $V(D)$）为沿 $G\to\mathrm{GL}(V)$ 前推所得的 $\Gal_F$-等变 $\mathrm{GL}(V)$-丛（相应地带 $\mathrm{GL}(V)$-结构的 $(\varphi,\Gamma)$-模）。我们仍把 $V(\mathcal{E})$ 看作向量丛而非 $\mathrm{GL}(V)$-丛，类似地把 $V(D)$ 看作 $(\varphi,\Gamma)$-模。对特征标 $\check\lambda\in X^*(T)^+_G$，记 $V^{\check\lambda}$ 为最高权 $\check\lambda$ 的不可约 $G$-表示，$L^{\check\lambda}$ 为 $T$ 的典范表示 $\check\lambda: T\to\mathbb{G}_m$。

下文关于 $\overline{\mathfrak{X}}^{\mathrm{naive}}_P$ 的定义以及关于 $\widetilde{\mathfrak{X}}^{\mathrm{naive}}_P$ 的定义，是 [BG02] 中定义的 $\Bun_P$ 相应紧化的类似物。

\begin{defn}{5.3.3}
叠 $\overline{\mathfrak{X}}^{\mathrm{naive}}_P$ 是 $\mathrm{Rig}_L$ 上的群胚纤维范畴（category fibered in groupoids），参数化三元组
{\small
\[
\overline{\mathfrak{X}}^{\mathrm{naive}}_P(\operatorname{Sp} A)=\Bigl\{\mathcal{E}_G\in\mathfrak{X}_G(\operatorname{Sp} A),\ \mathcal{E}_{S_M}\in\mathfrak{X}_{S_M}(\operatorname{Sp} A);\ \bigl(\kappa_{\check\lambda}: L^{\check\lambda}(\mathcal{E}_{S_M})\to V^{\check\lambda}(\mathcal{E}_G)\bigr)_{\check\lambda\in X^*(T)^+_G\cap X^*(T)_{G,P}}\Bigr\},
\]}
使得 $\kappa_{\check\lambda}$ 到 $\operatorname{Sp} A$ 上纤维的限制是单射、各 $\kappa_{\check\lambda}$ 的余核在 $\operatorname{Sp} A$ 上平坦，且满足下述 Plücker 关系：
\ysrc{226}
\begin{itemize}[leftmargin=2.2em,itemsep=0.3em]
\item 当 $\check\lambda=0$ 时，映射 $\kappa_0$ 是恒等。
\item 对 $\check\lambda,\check\mu\in X^*(T)^+_G\cap X^*(T)_{G,P}$，复合
\[
\begin{aligned}
L^{\check\lambda}(\mathcal{E}_{S_M})\otimes L^{\check\mu}(\mathcal{E}_{S_M})
&=L^{\check\lambda+\check\mu}(\mathcal{E}_{S_M})\xrightarrow{\ \kappa_{\check\lambda+\check\mu}\ }V^{\check\lambda+\check\mu}(\mathcal{E}_G)\\
&\longrightarrow(V^{\check\lambda}\otimes V^{\check\mu})(\mathcal{E}_G)=V^{\check\lambda}(\mathcal{E}_G)\otimes V^{\check\mu}(\mathcal{E}_G)
\end{aligned}
\]
与映射 $\kappa_{\check\lambda}\otimes\kappa_{\check\mu}$ 一致。
\end{itemize}
\label{thm:a4-5-3-3}
\end{defn}

$X_{\bar F}$ 上等变 $P$-丛的叠 $\mathfrak{X}_P$ 可与 $\overline{\mathfrak{X}}^{\mathrm{naive}}_P$ 的开子叠等同，即各 $\kappa_{\check\lambda}$ 的余核在 $\operatorname{Sp} A\times X_{\bar F}$ 上（而不只是在 $\operatorname{Sp} A$ 上）平坦之处。由定义显然，映射 $\mathfrak{X}_P\to\mathfrak{X}_G$ 与 $\mathfrak{X}_P\to\mathfrak{X}_M\to\mathfrak{X}_{S_M}$ 可延拓为态射
\[
\overline{\beta}_P: \overline{\mathfrak{X}}^{\mathrm{naive}}_P\longrightarrow\mathfrak{X}_G,\qquad
\overline{\alpha}_P: \overline{\mathfrak{X}}^{\mathrm{naive}}_P\longrightarrow\mathfrak{X}_{S_M},
\]
但 $\overline{\alpha}$ 不经过 $\mathfrak{X}_M$ 分解。此外，$P\mapsto\overline{\mathfrak{X}}^{\mathrm{naive}}_P$ 对抛物子群的包含 $P'\subseteq P$ 具有函子性（与投影映射 $\overline{\alpha}$、$\overline{\beta}$ 相容）。

\begin{defn}{5.3.4}
叠 $\widetilde{\mathfrak{X}}^{\mathrm{naive}}_P$ 是 $\mathrm{Rig}_L$ 上的群胚纤维范畴，参数化三元组
\[
\widetilde{\mathfrak{X}}^{\mathrm{naive}}_P(\operatorname{Sp} A)=\Bigl\{\mathcal{E}_G\in\mathfrak{X}_G(\operatorname{Sp} A),\ \mathcal{E}_M\in\mathfrak{X}_M(\operatorname{Sp} A);\ \bigl(\kappa_V: V^{U_P}(\mathcal{E}_M)\to V(\mathcal{E}_G)\bigr)_V\Bigr\},
\]
其中映射 $\kappa_V$ 以有限维 $G$-表示为下标，使得 $\kappa_V$ 到 $\operatorname{Sp} A$ 上纤维的限制是单射、各 $\kappa_V$ 的余核在 $\operatorname{Sp} A$ 上平坦，且满足下述版本的 Plücker 关系：
\begin{itemize}[leftmargin=2.2em,itemsep=0.3em]
\item 若 $V$ 是平凡表示，则映射 $\kappa_V$ 是恒等。
\item 对任意两个 $G$-表示 $V_1$、$V_2$，下图交换：
\[
\begin{tikzcd}[column sep=large,row sep=large]
(V_1^{U_P}\otimes V_2^{U_P})(\mathcal{E}_M) \arrow[r,"\kappa_{V_1}\otimes\kappa_{V_2}"] \arrow[d] & (V_1\otimes V_2)(\mathcal{E}_G) \arrow[d,"="]\\
(V_1\otimes V_2)^{U_P}(\mathcal{E}_M) \arrow[r,"\kappa_{V_1\otimes V_2}"'] & (V_1\otimes V_2)(\mathcal{E}_G).
\end{tikzcd}
\]
\item 对 $G$-表示的态射 $V_1\to V_2$，下图交换：
\[
\begin{tikzcd}[column sep=large,row sep=large]
V_1^{U_P}(\mathcal{E}_M) \arrow[r,"\kappa_{V_1}"] \arrow[d] & V_1(\mathcal{E}_G) \arrow[d]\\
V_2^{U_P}(\mathcal{E}_M) \arrow[r,"\kappa_{V_2}"'] & V_2(\mathcal{E}_G).
\end{tikzcd}
\]
\end{itemize}
\label{thm:a4-5-3-4}
\end{defn}

同样，$\mathfrak{X}_P$ 可与 $\widetilde{\mathfrak{X}}^{\mathrm{naive}}_P$ 中使各 $\kappa_V$ 的余核平坦（在 $\operatorname{Sp} A\times X_{\bar F}$ 上，而不只是 $A$-平坦）的开子叠等同。同样由定义显然，
